\documentclass[11pt]{article}
\usepackage[utf8]{inputenc}
\usepackage{amsmath, amssymb}
\usepackage{geometry}
\geometry{a4paper, margin=1in}

\title{Detailed Explanation of Slide 9: GIV Supply and Demand Elasticities}
\author{Granular Instrumental Variables (Gabaix and Koijen, 2022)}
\date{}

\begin{document}

\maketitle

\section{Introduction}
Slide 9 of the presentation demonstrates how the Granular Instrumental Variable ($z_t$) can be used to identify both the supply elasticity ($\psi$) and the demand elasticity ($\phi^d$). This document provides the underlying derivations and the statistical assumptions required for these results.

\section{The Core Model and the GIV}
Recall the system of equations from slide 5:
\begin{align}
p_t &= \psi y_{St} + \epsilon_t \label{eq:supply} \tag{38} \\
y_{it} &= \phi^d p_t + \eta_t + u_{it} \label{eq:demand_firm} \tag{39}
\end{align}
where $y_{St} = \sum S_i y_{it}$ is the size-weighted aggregate outcome. From slide 8, we define the Granular Instrumental Variable ($z_t$) as the difference between size-weighted and equal-weighted aggregate outcomes:
\begin{equation}
z_t = y_{St} - y_{Et} = u_{St} - u_{Et} \label{eq:giv_def} \tag{42}
\end{equation}
where $u_{St} = \sum S_i u_{it}$ and $u_{Et} = \frac{1}{N} \sum u_{it}$. 

\section{The Exogeneity Condition}
The fundamental assumption of the GIV approach is that the idiosyncratic shocks $u_{it}$ are independent of the aggregate shocks $\epsilon_t$ and $\eta_t$. Since $z_t$ is a linear combination of only idiosyncratic shocks, it follows that:
\begin{equation}
\mathbb{E}[z_t \epsilon_t] = 0 \quad \text{and} \quad \mathbb{E}[z_t \eta_t] = 0 \label{eq:exogeneity}
\end{equation}
This exogeneity allows $z_t$ to serve as a valid instrument.

\section{Identifying Supply Elasticity ($\psi$)}
To find $\psi$, we use $z_t$ as an instrument for $y_{St}$ in the supply equation (\ref{eq:supply}). We take the expectation of the product of the supply equation and the instrument:
\begin{align*}
\mathbb{E}[p_t z_t] &= \mathbb{E}[(\psi y_{St} + \epsilon_t) z_t] \\
\mathbb{E}[p_t z_t] &= \psi \mathbb{E}[y_{St} z_t] + \mathbb{E}[\epsilon_t z_t]
\end{align*}
By the exogeneity condition in (\ref{eq:exogeneity}), $\mathbb{E}[\epsilon_t z_t] = 0$. Thus:
\begin{equation}
\psi = \frac{\mathbb{E}[p_t z_t]}{\mathbb{E}[y_{St} z_t]}
\end{equation}
This is the standard IV formula where $z_t$ instruments for $y_{St}$.

\section{Identifying Demand Elasticity ($\phi^d$)}
To find $\phi^d$, we use the equal-weighted aggregate demand equation derived on slide 8:
\begin{equation}
y_{Et} = \phi^d p_t + \eta_t + u_{Et} \tag{41}
\end{equation}
We aim to show that $\mathbb{E}[(y_{Et} - \phi^d p_t) z_t] = 0$. Substituting the demand equation:
\begin{equation}
\mathbb{E}[(y_{Et} - \phi^d p_t) z_t] = \mathbb{E}[(\eta_t + u_{Et}) z_t] = \mathbb{E}[\eta_t z_t] + \mathbb{E}[u_{Et} z_t]
\end{equation}
We already know $\mathbb{E}[\eta_t z_t] = 0$. Now we examine $\mathbb{E}[u_{Et} z_t]$. Under the assumption that $u_{it}$ are i.i.d. with variance $\sigma_u^2$:
\begin{align*}
\mathbb{E}[u_{Et} z_t] &= \mathbb{E}[u_{Et} (u_{St} - u_{Et})] \\
&= \mathbb{E}[u_{Et} u_{St}] - \mathbb{E}[u_{Et}^2] \\
&= \mathbb{E} \left[ \left( \frac{1}{N} \sum_{i} u_{it} \right) \left( \sum_{j} S_j u_{jt} \right) \right] - \mathbb{E} \left[ \left( \frac{1}{N} \sum_{i} u_{it} \right)^2 \right] \\
&= \frac{1}{N} \sum_{i} S_i \sigma_u^2 - \frac{1}{N^2} \sum_{i} \sigma_u^2 \\
&= \frac{1}{N} \sigma_u^2 (1) - \frac{1}{N} \sigma_u^2 = 0
\end{align*}
Since $\mathbb{E}[u_{Et} z_t] = 0$, it follows that $\mathbb{E}[(y_{Et} - \phi^d p_t) z_t] = 0$. Rearranging:
\begin{align*}
\mathbb{E}[y_{Et} z_t] &= \phi^d \mathbb{E}[p_t z_t] \\
\phi^d &= \frac{\mathbb{E}[y_{Et} z_t]}{\mathbb{E}[p_t z_t]}
\end{align*}
This identifies the demand elasticity using the same granular instrument $z_t$.

\end{document}
